\documentclass[11pt]{article}
\usepackage[T1]{fontenc}
\usepackage{lmodern}
\usepackage[margin=1in]{geometry}
\usepackage{amsmath,amssymb,amsthm}
\usepackage{microtype}
\usepackage[hidelinks]{hyperref}
\hypersetup{pdftitle={A bijection for 312/321-avoiding permutations whose cubes avoid 2143},
  pdfauthor={Xiaojian Zeng}}
\newtheorem{theorem}{Theorem}
\newtheorem{lemma}[theorem]{Lemma}
\newtheorem{corollary}[theorem]{Corollary}
\newcommand{\Av}{\operatorname{Av}}
\newcommand{\eps}{\varepsilon}
\newcommand{\C}{\mathcal C}
\newcommand{\B}{\mathcal B}

\title{A bijection for $312/321$-avoiding permutations\\
whose cubes avoid $2143$}
\author{Xiaojian Zeng\\[3pt]
\small PBC School of Finance, Tsinghua University, 100190\\
\small\href{mailto:xiaojianzeng@foxmail.com}{xiaojianzeng@foxmail.com}}
\date{}

\begin{document}
\maketitle

\begin{abstract}
We prove the cube-enumeration conjecture in Section~5 of Archer and
Bourne's \emph{Pattern avoidance in compositions and powers of
permutations}. Their cyclic-block encoding restricts to a bijection
between permutations avoiding $312$ and $321$ whose cubes avoid $2143$
and ordered compositions with at most one part outside $\{1,3\}$.
The inverse reads successive positions of the minimum in the original
permutation. The restriction follows because every cubed block avoids
$2143$, while two blocks containing inversions create $2143=21\oplus21$.
We obtain the resulting rational generating function, including the empty
objects. No condition on the square of the permutation is imposed.
\end{abstract}

\section{Statement and context}

For $n\geq0$, let $S_n$ be the permutations of $[n]=\{1,\ldots,n\}$,
with one empty permutation in $S_0$. We use one-line notation. A
permutation contains a \emph{classical pattern} $\tau$ if a subsequence
of its entries has the same relative order as $\tau$; the selected
positions need not be consecutive. Write $\Av_n(312,321)$ for the
permutations in $S_n$ avoiding both patterns. Powers are compositional:
$\pi^3(i)=\pi(\pi(\pi(i)))$.

Pattern restrictions on permutation powers extend the strong avoidance
studied by B\'ona and Smith~\cite{BonaSmith2019}, where both a permutation
and its square avoid the same pattern. Archer and Geary
\cite{ArcherGeary2024} study chains that permit different restrictions
on different powers. Archer and Bourne~\cite{ArcherBourne2026} encode
$\Av_n(312,321)$ by compositions and use this encoding to enumerate
square-pattern restrictions. In their Section~5 they conjecture the
cube enumeration proved here.

A \emph{composition} of $n$ is an ordered list of positive integers
summing to $n$; only the empty list composes zero. Define
\begin{align*}
\B_n&=\{\pi\in\Av_n(312,321):\pi^3\text{ avoids }2143\},\\
\C_n&=\{(d_1,\ldots,d_m):d_i\geq1,\ \textstyle\sum_i d_i=n,
          \ |\{i:d_i\notin\{1,3\}\}|\leq1\}.
\end{align*}
An exceptional part may occur in any position and carries no additional
mark. There is no restriction on $\pi^2$ in the definition of $\B_n$.

For $\alpha\in S_a$ and $\beta\in S_b$, their \emph{direct sum}
$\alpha\oplus\beta$ is the concatenation of $\alpha$ and the entries of
$\beta$ increased by $a$. Set
\[
\eps_1=(1),\qquad \eps_d=(2,3,\ldots,d,1)\quad(d\geq2),\qquad
F(d_1,\ldots,d_m)=\bigoplus_{i=1}^{m}\eps_{d_i}.
\]
The empty direct sum is the empty permutation.

\begin{theorem}\label{thm:main}
For every $n\geq0$, $F$ restricts to a size-preserving bijection
$\C_n\longrightarrow\B_n$. Its inverse repeatedly records the position
$d$ of the minimum, removes the prefix $\eps_d$, and subtracts $d$
from every remaining entry. More precisely, for every finite list of
positive integers,
\begin{equation}\label{eq:criterion}
F(d_1,\ldots,d_m)^3\text{ avoids }2143
\quad\Longleftrightarrow\quad
|\{i:d_i\notin\{1,3\}\}|\leq1.
\end{equation}
\end{theorem}

The unrestricted encoding is Archer--Bourne's Lemma~3.1
\cite{ArcherBourne2026}. Compatibility of powers with direct sums is
also known~\cite[Lemma~2.1]{ArcherGeary2024}. The contribution here is
the exact restriction~\eqref{eq:criterion} and the resulting enumeration.
The closest higher-power comparison is Archer--Geary's
Theorem~5.2~\cite{ArcherGeary2024}: for $231/321$-avoiders whose
$k$th powers avoid $132$, every part after the first must divide $k$.
Here the target is $2143$, and an exceptional part is allowed anywhere.
The proof below supplies all the required structural arguments.

\section{The restricted bijection}

We first recall the block encoding, keeping its inverse explicit.

\begin{lemma}[Archer--Bourne]\label{lem:encoding}
The map $F$ is a bijection from all compositions of $n$ to
$\Av_n(312,321)$, for every $n\geq0$.
\end{lemma}

\begin{proof}
In a block $\eps_d$, each entry has at most one smaller entry to its
right. Later direct-sum blocks have larger values, so the same assertion
holds in $F(d_1,\ldots,d_m)$. The first entry of either a $312$ or a
$321$ occurrence would have two smaller entries to its right. Hence
every image of $F$ avoids both patterns.

Conversely, let $\pi\in\Av_n(312,321)$ be nonempty, and let $d$ be the
position of $1$. If $i<d<j$ and $\pi_i>\pi_j$, then
$(\pi_i,1,\pi_j)$ is a $312$ occurrence. Thus every value before $1$
is smaller than every value after it, and the prefix of length $d$
has value set $[d]$. If $i<j<d$ and $\pi_i>\pi_j$, then
$(\pi_i,\pi_j,1)$ is a $321$ occurrence. The prefix is therefore
exactly $\eps_d$. Subtracting $d$ from the suffix yields
$\rho\in\Av_{n-d}(312,321)$, and $\pi=\eps_d\oplus\rho$.
This argument includes $d=1$ and $d=n$.

Define $G(\varnothing)=()$ and obtain $G(\pi)$ by placing $d$ before
the list $G(\rho)$. The recursion terminates since each step removes
at least one entry. In $F(d_1,\ldots,d_m)$, the minimum occurs at the
end of the first block, in position $d_1$; removing this block leaves
the encoding of the remaining parts after subtracting $d_1$.
Induction on the number of parts gives $G(F(\mathbf d))=\mathbf d$.
Conversely, the forced-prefix decomposition and induction on $n$ give
$F(G(\pi))=\eps_d\oplus F(G(\rho))=\pi$. Both inductions start
with the empty objects. The maps preserve the sum of the block lengths.
\end{proof}

\begin{lemma}\label{lem:cubes}
Let $q_d=\eps_d^3$. Then $q_d$ avoids $2143$ for every $d\geq1$,
and $q_d$ has an inversion if and only if $d\notin\{1,3\}$.
Moreover,
\begin{equation}\label{eq:power-sum}
F(d_1,\ldots,d_m)^3=q_{d_1}\oplus\cdots\oplus q_{d_m}.
\end{equation}
\end{lemma}

\begin{proof}
On $[d]$, the block $\eps_d$ adds one modulo $d$, with residues
represented by $1,\ldots,d$. Thus
\[
q_d(i)=1+((i-1+3)\bmod d),
\]
and in one-line notation
\begin{equation}\label{eq:blocks}
q_1=(1),\quad q_2=(2,1),\quad q_3=(1,2,3),\quad
q_d=(4,5,\ldots,d,1,2,3)\quad(d\geq4).
\end{equation}
The cases $d=1,3$ are increasing. The block $q_2$ has an inversion,
and for $d\geq4$ the consecutive entries $d,1$ give one.

For $d=2$ and $d\geq4$, the block $q_d$ in~\eqref{eq:blocks} is a
concatenation of two increasing runs. Any subsequence takes entries from the first
run followed by entries from the second, and so has at most one descent
between successive selected entries. Since $2143$ has two such
descents, it cannot occur. The increasing cases also avoid $2143$.

Finally, each direct-sum block occupies the same interval of positions
and values, which is invariant under $F$. If its offset is $s$, then
$F(s+i)=s+\eps_d(i)$ inside that block. Applying this identity three
times proves~\eqref{eq:power-sum}.
\end{proof}

\begin{lemma}\label{lem:two-inversions}
Suppose $\sigma_1,\ldots,\sigma_m$ are nonempty permutations, each
avoiding $2143$. Their direct sum avoids $2143$ if and only if at most
one summand has an inversion. The assertion also holds for the empty
direct sum.
\end{lemma}

\begin{proof}
Every inversion in a direct sum lies within a summand: an earlier
summand has smaller positions and smaller values than a later one.
If two summands have inversions, choose positions $a<b$ in the earlier
one and $c<e$ in the later one. Writing $\sigma$ for the direct sum,
we have
\[
a<b<c<e,\qquad
\sigma(b)<\sigma(a)<\sigma(e)<\sigma(c).
\]
These four entries form $2143$.

Conversely, an occurrence of $2143$ at $a<b<c<e$ supplies the
inversions $(a,b)$ and $(c,e)$. Each pair lies in a summand. If they
lay in the same summand, all four selected entries would lie there,
contradicting that summand's avoidance of $2143$. An occurrence therefore
requires two distinct summands with inversions. The empty sum contains
neither an inversion nor a four-entry pattern.
\end{proof}

\begin{proof}[Proof of Theorem~\ref{thm:main}]
By Lemma~\ref{lem:cubes}, the summands of $F(\mathbf d)^3$ all avoid
$2143$, and precisely the parts outside $\{1,3\}$ yield summands
with inversions. Lemma~\ref{lem:two-inversions} now gives
\eqref{eq:criterion}. Consequently $F$ maps $\C_n$ into $\B_n$,
and its inverse $G$ from Lemma~\ref{lem:encoding} maps $\B_n$ into
$\C_n$. The two inverse identities restrict to these sets, including
the empty objects.
\end{proof}

The inverse must act on the original permutation: the compositions
$(3)$ and $(1,1,1)$ encode different permutations with equal cubes.
Also, internal avoidance is essential in Lemma~\ref{lem:two-inversions};
the single summand $2143$ would violate its conclusion if that hypothesis
were removed.

\section{Enumeration and finite checks}

\begin{corollary}\label{cor:gf}
Writing $b_n=|\B_n|$, one has $b_n=|\C_n|$ for every $n\geq0$ and
the formal power-series identity
\begin{equation}\label{eq:gf}
\sum_{n\geq0}b_nz^n
=\frac{1}{1-z-z^3}
 +\frac{z^2+z^4/(1-z)}{(1-z-z^3)^2}.
\end{equation}
\end{corollary}

\begin{proof}
Let $A(z)$ count sequences of parts in $\{1,3\}$ and let $E(z)$
count a single exceptional part. Then
\[
A(z)=\sum_{r\geq0}(z+z^3)^r=\frac{1}{1-z-z^3},\qquad
E(z)=\sum_{\substack{d\geq1\\d\notin\{1,3\}}}z^d
    =z^2+\frac{z^4}{1-z}.
\]
An admissible composition either has no exceptional part or has a
unique decomposition into an ordinary prefix, an exceptional part,
and an ordinary suffix. These disjoint cases contribute $A$ and
$AEA$, respectively. Theorem~\ref{thm:main} therefore gives
\eqref{eq:gf}. All operations take place in $\mathbb Z[[z]]$:
positive part sizes make each coefficient sum finite, and every
denominator has constant term one. The constant coefficient is $1$.
\end{proof}

The first values are
\[
(b_0,b_1,\ldots,b_9)=(1,1,2,4,7,13,23,39,66,110).
\]
As finite diagnostics, a separately written standard-library Python
verifier examined all $409{,}114$ permutations of sizes $0$ through $9$.
It tested input avoidance by direct subsequence ranks, formed each cube
by three literal applications, and compared the entire resulting sets
with the unrestricted and restricted block images. Both inverse
identities were checked. It also tested Lemma~\ref{lem:two-inversions}
on all $1{,}057$ lists of length at most two from the $32$ nonempty
$2143$-avoiders of sizes at most four, and compared the coefficients
of~\eqref{eq:gf} with a last-part composition recurrence through
degree $100$. All these finite comparisons agreed.

Using exact predicates and a separate implementation to recheck
computed objects has a methodological parallel in the certificate
architecture in Section~6 of Zeng, Liu, Ratnavelu, and Ong
\cite{ZengEtAl2026}. Here the separate verifier checks for
encoding, iteration, and pattern-detection errors. The computations
are supplementary finite observations; the lemmas and the unique
exceptional-part decomposition prove the assertion for every $n$.

\paragraph{Reproducibility and assistance.}
The accompanying computational supplement includes the two exact
Python verification programs, their recorded outputs, and run
instructions. OpenAI Codex assisted with proof development, literature
search, verification code, and manuscript preparation. A separate
model-context pass checked the argument and wrote the second verifier;
this is neither external peer review nor statistically independent
evidence. No proof assistant was used. The result resolves the specific
cube conjecture stated above; it does not claim a classification for
arbitrary powers or target patterns. The comparison with prior work is
limited to the cited sources and targeted searches through
12 September 2026.

\bibliographystyle{plainurl}
\bibliography{references}
\end{document}
